% ===================================================================
%  TABELO N-RO 16 — La granda magazeno. — La bazaro. — La ludiloj.
%  Traduction 2026 d'apres texto/io, controlee sur texto/fr et
%  texto/es. Consigne : voir texto/eo/10-tabelo-01.tex.
%
%  Le tableau d'astronomie (bloc 15-1) porte des renvois a LETTRES,
%  (a) a (f), graves sur le tableau lui-meme : ils se rangent sous le
%  numero 85, comme le dit gravuri/literi.json.
%
%  L'IDO SAUTE LA LETTRE (d) DE L'ENUMERATION des cartes, que le
%  fac-simile grave pourtant. La colonne esperantiste la rend, comme
%  les autres traductions : c'est le troisieme des trois ecarts
%  declares dans outils/renvoji.py.
% ===================================================================

%%K t16-apar-1 apar
\VUsaut{4.0mm}
\VUtitre{15.0pt}{60}{TABELO N\textsuperscript{o} 16}
\VUsaut{2.0mm}
\VUpk{13.2pt}{40}{La granda magazeno. --- La bazaro.}
\VUpk{13.2pt}{40}{La ludiloj.}
\VUsaut{3.4mm}

%%K t16-01-1 p
1. --- La nova jaro alproksimiĝas. La grandaj magazenoj preparis siajn
elmetejojn de \VUgras{ludiloj} \textsuperscript{(1)} kaj novjaraj donacoj.

%%K t16-01-2 p
Mi petis de mia avo, ke li konduku min en la bazaron por vidi tiun belan
ekspoziciejon, kaj li konsentis tion.

%%K t16-02-1 p
2. --- Unue ni haltis antaŭ belaj armoŝildoj, kiuj enhavas \VUgras{fusilon},
\VUgras{kepon}, \VUgras{sabron}, \VUgras{spadzonon} kaj paron da
\VUgras{epoletoj}, aŭ \VUgras{kaskon}, \VUgras{kirason} \textsuperscript{(3)}
kaj \VUgras{pistolon} \textsuperscript{(4)}. Tio ĉio interesis min multe pli ol
la \VUgras{lumprojekciiloj} \textsuperscript{(2)}.

%%K t16-tit-1 sub
\VUsaut{3.0mm}
\VUpk{10.2pt}{40}{Belaj vidaĵoj.}
\VUsaut{2.6mm}

%%K t16-03-1 p
3. --- En la sama fako estis \VUgras{gardostaranto} \textsuperscript{(5)} en sia
\VUgras{gardobudo}; li ŝajnis gardostari antaŭ tuta menaĝerio kartona. Kiom
multaj bestoj estis tie: \VUgras{Papago} \textsuperscript{(6)} sur sia
sidstango, \VUgras{rinocero} \textsuperscript{(7)} kun korno sur la nazo,
\VUgras{boaco} \textsuperscript{(8)}, \VUgras{struto} \textsuperscript{(9)},
\VUgras{simio} \textsuperscript{(10)} krakanta juglandon, \VUgras{vulpo}
\textsuperscript{(11)} forportanta kokinon, \VUgras{krokodilo}
\textsuperscript{(12)} gapanta per enormaj makzeloj, \VUgras{kamelo}
\textsuperscript{(13)}, eleganta \VUgras{leporhundo} \textsuperscript{(14)},
\VUgras{pantero} \textsuperscript{(15)}, poste du bestoj kiujn mi ne konis kaj
kiuj estas, laŭ diro, \VUgras{mustelo} \textsuperscript{(16)} kaj \VUgras{hieno}
\textsuperscript{(17)}. Ankaŭ tie estis \VUgras{leopardo}
\textsuperscript{(18)} kaj \VUgras{lupo} \textsuperscript{(21)}; tute proksime
estis ĉarmaj \VUgras{ratetoj} \textsuperscript{(19)} kaj etegaj \VUgras{musetoj}
\textsuperscript{(20)} kaŭĉukaj. Laste \VUgras{hipopotamo}
\textsuperscript{(22)} kaj \VUgras{dromedaro} \textsuperscript{(23)} ĝiba finis
la kolektaĵon. Mi estus restinta tie dum pli multaj horoj se ne estus estinta
apude belega kampejo plena de \VUgras{plumbaj soldatoj} \textsuperscript{(29)}
kiu altiris mian atenton.

%%K t16-tit-2 sub
\VUsaut{4.0mm}
\VUpk{10.2pt}{40}{Soldatoj kaj milito.}
\VUsaut{2.8mm}

%%K t16-04-1 p
4. --- Mia avo, kiu estas \VUgras{invalido} \textsuperscript{(24)} kaj kiu devis
forlasi la armeon pro \VUgras{vundo} kiu igis lin akiri la \VUgras{militistan
medalon}, tre ŝatas rakonti al mi la \VUgras{kampanjojn} en kiuj li
partoprenis. Li estis feliĉa klarigante la diversajn movadojn de la malgranda
armeo miniatura. Grupo de veraj soldatoj vizitanta la bazaron: \VUgras{genia
soldato} \textsuperscript{(25)}, \VUgras{alpa ĉasisto} \textsuperscript{(26)}
kapkovrita per \VUgras{bereto}, \VUgras{majorserĝento} \textsuperscript{(27)} de
infanterio kaj \VUgras{artileria serĝento} \textsuperscript{(28)} kun ora
\VUgras{galono} sur maniko, haltis apud ni por aŭskulti la klarigojn.

%%K t16-05-1 p
5. --- Mia avo, tute fiera, ĉar li havis tiel brilan aŭdantaron, improvizis
kompletan batalplanon.

%%K t16-05-2 p
La fortikigita urbo, kun siaj \VUgras{remparoj} kaj \VUgras{fosaĵoj}, estis
sieĝata de la \VUgras{armeo malamika}, kiu, post longtempa \VUgras{bombardado},
estis preta por \VUgras{ataki}. Sed la \VUgras{sieĝatoj} devigitaj pro la
malsato, estis provintaj \VUgras{elataki} al la \VUgras{kampejo} de la
\VUgras{sieĝantoj}. Post repuŝi la \VUgras{atakon} per obstina \VUgras{batalo}
ĉirkaŭ la \VUgras{tendaro}, la \VUgras{venkintaj} sieĝantoj lanĉis siajn
\VUgras{eskadronojn} por persekuti la \VUgras{venkitajn} atakintojn. Ĉi tiuj
\VUgras{retiriĝis} kovrante la \VUgras{batalkampon} per \VUgras{mortintoj} kaj
\VUgras{vunditoj}. \VUgras{Brankardistoj} forportis tiujn lastajn en la
\VUgras{ambulancon} por flegigi ilin de la \VUgras{kirurgo}.

%%K t16-05-3 p
Malgraŭ la heroa rezisto de kelkaj \VUgras{batalionoj} kiuj estis formintaj
\VUgras{kvadraton}, la \VUgras{venkiĝo} estis kompleta, la cetero de la armeo
\VUgras{forkuris}, lasante multajn \VUgras{kaptitojn}. La malamiko iris
finfari sian \VUgras{venkon} okupante la urbon. Eble tio estos la fino de la
kampanjo kaj, post \VUgras{armistico}, estos eble subskribi
\VUgras{packontrakton}.

%%K t16-tit-3 sub
\VUsaut{4.4mm}
\VUpk{10.2pt}{40}{Donacaĵoj por novjaro.}
\VUsaut{2.8mm}

%%K t16-06-1 p
6. --- La junaj soldatoj, kiuj ridis aŭskultante la pentrindan rakonton de la
veterano petis de li kelkajn aliajn klarigojn. Pri mi, mi ĉirkaŭiris la
vendejon por rigardi belan \VUgras{arbaleston} \textsuperscript{(32)} kaj
\VUgras{arkon} \textsuperscript{(33)} kun ĝia \VUgras{sago}
\textsuperscript{(31)}. Tie mi trovis alian kolektaĵon da bestoj: du
\VUgras{birdoj kantantaj} \textsuperscript{(34)}: aŭtomata \VUgras{najtingalo}
kaj \VUgras{silvio} kantanta per tuta gorĝo, kaj serio da strangaj bestoj:
\VUgras{vesperto} \textsuperscript{(35)}, \VUgras{boao} \textsuperscript{(36)},
\VUgras{testudo} \textsuperscript{(37)}, \VUgras{foko} \textsuperscript{(38)},
\VUgras{baleno} \textsuperscript{(39)} kaj \VUgras{ŝarko}
\textsuperscript{(40)} el orita ledo.

%%K t16-07-1 p
7. --- Mi ne longe haltis por kontempli ilin, ĉar mi ekvidis la revatan:
belegan \VUgras{marionetteatron} \textsuperscript{(41)} kun belegaj
\VUgras{dekoraĵoj} kaj raviga \VUgras{kurteno} ruĝa. Sur la \VUgras{sceno}, du
aktoroj ŝajnis ludi \VUgras{rolon} en tre interesa \VUgras{teatraĵo}, ĉar la
malgrandaj \VUgras{spektantoj} kartonaj instalitaj en la loĝioj aŭ sidantaj sur
la \VUgras{brakseĝoj} kaj en la \VUgras{partero} malantaŭ la
\VUgras{orkestrestro}, aplaŭdis vigle.

%%K t16-07-2 p
Bedaŭrinde, tiu teatro kostis tre kare.

%%K t16-08-1 p
8. --- Cetere, elekti la ludilojn estis malfacile, ĉar en la vitrinoj estis
amasigitaj ĉiaspecaj ludiloj: \VUgras{Arlekeno} \textsuperscript{(42)},
\VUgras{Pulĉinelo} \textsuperscript{(43)}, \VUgras{Pierrot}
\textsuperscript{(44)}, \VUgras{gufoj} \textsuperscript{(45)} svingantaj la
flugilojn, fajfantaj \VUgras{merloj} \textsuperscript{(46)}, bojantaj
\VUgras{pudeloj}, \VUgras{fonografilo} \textsuperscript{(47)} kaj
\VUgras{ludiloj de pacienco}. Unu el tiuj ludiloj tre amuzis min: ĝi
prezentis \VUgras{maran batalon} \textsuperscript{(48)} en kiu \VUgras{ŝipo}
estis atakata de \VUgras{piratoj} apud lando plantita per \VUgras{kokosarboj}.

%%K t16-noto-1 noto
\VUnotes{143.2mm}{%
\textit{(*) Vidu la tabelon N\textsuperscript{o} 6.}%
}

%%K t16-09-1 p
9. --- Mi povis tre facile kontempli tiujn ĉiujn mirindaĵojn, ĉar nur malmultaj
personoj ĉeestis en la magazeno, apenaŭ kelkaj flanistoj, \VUgras{dragono}
\textsuperscript{(49)} kaj \VUgras{enkasigisto} \textsuperscript{(50)} kiu
prezentis \VUgras{traton} ĉe la ĉefa \VUgras{kaso} \textsuperscript{(51)}.

%%K t16-10-1 p
10. --- Mi demandis mian avon pri la uzo de la libroj metitaj sur bretojn
malantaŭ la \VUgras{kasistino}; li respondis al mi ke ili estas
\VUgras{kontolibroj}.

%%K t16-tit-4 sub
\VUsaut{3.6mm}
\VUpk{10.2pt}{40}{Gapado ĉirkaŭ butiko.}
\VUsaut{2.8mm}

%%K t16-11-1 p
11. --- Forlasante la fakon de la ludiloj, ni iris al la selfarejo doni rigardon
al la \VUgras{seloj} \textsuperscript{(53)}, \VUgras{piedingoj}
\textsuperscript{(54)}, \VUgras{bridoj} \textsuperscript{(55)},
\VUgras{mordaĵoj} \textsuperscript{(56)} kaj \VUgras{rajdvipoj}. Dum la
\VUgras{fakestro} \textsuperscript{(57)} donis kelkajn informojn al mia avo, mi
iris rigardi la elmetejon de la \VUgras{porcelanoj} kaj \VUgras{kristaloj}
\textsuperscript{(58)} kie estas belaj \VUgras{salatujoj} kaj delikataj
\VUgras{tekruĉoj} \textsuperscript{(59)}.

%%K t16-12-1 p
12. --- Por iri sur la unuan etaĝon, ni pasis malantaŭ la vitrino de la
\VUgras{korsetoj} \textsuperscript{(60)}, apud la ŝtrumpvendejo, kie estis
elmetitaj tre belaj \VUgras{pantalonoj} \textsuperscript{(63)}. Proksime,
\VUgras{komizino} \textsuperscript{(61)} igis elekti de \VUgras{aĉetantino}
\textsuperscript{(62)} belajn silkajn \VUgras{teksaĵojn} \textsuperscript{(64)}.

%%K t16-12-2 p
Supreniranta per la ŝtuparo, mi rimarkis ke la magazeno estas dekorita per
japanaj \VUgras{lanternoj} \textsuperscript{(65)}, \VUgras{ombreloj}
\textsuperscript{(66)} kaj grandegaj \VUgras{palmarboj} \textsuperscript{(70)}.

%%K t16-13-1 p
13. --- Sur la unua etaĝo, ne estis multo vidinda; nur mebloj (*),
\VUgras{ferkestoj} \textsuperscript{(67)}, \VUgras{komodoj}
\textsuperscript{(68)}, \VUgras{infanveturiloj} \textsuperscript{(69)}. Sed la
ensembla vido de la magazeno estis tre \VUgras{amuza}.

%%K t16-14-1 p
14. --- Ĉe niaj piedoj, mi vidis la muzikilojn: \VUgras{harpoj}
\textsuperscript{(71)}, \VUgras{gitaroj} \textsuperscript{(72)},
\VUgras{mandolinoj} \textsuperscript{(73)}, \VUgras{valvkornetoj}
\textsuperscript{(74)}, \VUgras{klarionetoj} \textsuperscript{(75)}, violonoj
kun ilia \VUgras{arĉo} \textsuperscript{(76)} kaj eĉ \VUgras{tamburego}
\textsuperscript{(77)}. Iom pli malproksime, ĉe la paperfako,
\VUgras{studento} \textsuperscript{(78)} marĉandis de la \VUgras{komizo}
\textsuperscript{(79)} kajerujon. Apud ili, mi distingis belan
\VUgras{alfabetlibron} \textsuperscript{(80)}.

%%K t16-14-2 p
Meze de la magazeno estis starigita alta \VUgras{Kristnaskarbo}
\textsuperscript{(81)} tute provizita per kandeloj, aĵetoj kaj ĉiaspecaj
ludiloj: \VUgras{lignaj ĉevaloj} \textsuperscript{(82)}, \VUgras{pupoj}
\textsuperscript{(83)}, k. c..

%%K t16-14-3 p
La \VUgras{parfumvendejo} \textsuperscript{(84)} kun siaj \VUgras{dentifricoj}
kaj diversaj \VUgras{parfumoj} disvastigis tre bonan odoron kiu elspiris al ni.

%%K t16-tit-5 sub
\VUsaut{3.4mm}
\VUpk{10.2pt}{40}{Ni aĉetas ion.}
\VUsaut{2.8mm}

%%K t16-15-1 p
15. --- Ĉar mia avo komencis juĝi tro longa la tempon, ni malsupreniris per la
granda ŝtuparo. Li iom haltis antaŭ \VUgras{astronomia tabelo}
\textsuperscript{(85)} prezentanta, li diris al mi, \VUgras{kometojn}
\textsuperscript{(a)}, \VUgras{lun- kaj suneklipsojn} \textsuperscript{(b)},
ventrozon kun la \VUgras{kvar punktoj kardinalaj} \textsuperscript{(c)}
\VUgras{(Nordo, Sudo, Oriento kaj Okcidento)}, mapon de la du
\VUgras{duonsferoj} \textsuperscript{(d)}, la \VUgras{planedojn}
\textsuperscript{(e)} kaj la \VUgras{sunsistemon} \textsuperscript{(f)}. Mi
komprenis nenion en tio ĉio, kaj mi estis multe pli interesata de la galonita
livreo de \VUgras{knabo liveristo} \textsuperscript{(86)} kiu pasis, kaj de la
belaj \VUgras{ventumiloj} \textsuperscript{(87)} kiuj estis apud mi, apud la
\VUgras{monujoj} \textsuperscript{(88)} kaj la \VUgras{paperujoj}
\textsuperscript{(89)}.

%%K t16-16-1 p
16. --- Mi ankaŭ admiris sur la elmetejo \VUgras{plumojn}
\textsuperscript{(90)}, kaj \VUgras{zonbukojn} \textsuperscript{(91)}, la belajn
\VUgras{florojn artefaritajn} \textsuperscript{(94)} gracie aranĝitajn en
korboj. La \VUgras{violoj}, \VUgras{ĝilfloroj}, \VUgras{hiacintoj}
\textsuperscript{(95)}, \VUgras{geranioj} \textsuperscript{(96)},
\VUgras{lilioj} \textsuperscript{(97)} kaj \VUgras{kapucenoj}
\textsuperscript{(98)} estis tre bone kopiitaj.

%%K t16-tit-6 sub
\VUsaut{4.4mm}
\VUpk{10.2pt}{40}{Donacaĵo de mia avo.}
\VUsaut{2.8mm}

%%K t16-17-1 p
17. --- Mi petis mian avon aĉeti por mi unu el la \VUgras{dentobrosoj}
\textsuperscript{(92)} kiuj estis en skatolo apud \VUgras{harpingloj}
\textsuperscript{(93)}. Li konsentis kaj mi iris mem pagi mian aĉetaĵon ĉe la
kaso, apud kiu oni preparis gravan \VUgras{ekspedaĵon} \textsuperscript{(100)}
por \VUgras{kliento} \textsuperscript{(99)}.

%%K t16-18-1 p
18. --- Subite malpeza tumulto okazis inter la personoj haltintaj apud la artaj
\VUgras{bronzaĵoj} \textsuperscript{(101)}. \VUgras{Ŝtelisto}
\textsuperscript{(102)} estis ĵus kaptita, en sia delikto, de \VUgras{inspektisto}
\textsuperscript{(103)}, dum li provis prirabi iun. Oni zorgis venigi policiston
por \VUgras{aresti} lin, kaj ĝuste kiam ni eliris, la deliktinto estis
kondukata en la malliberejon.
\VUsaut{6.0mm}
\VUfilet{22mm}
